% !TeX program = lualatex
% START HERE: this document is both the template showcase and its user guide.
% For a clean manuscript, copy starter.tex to main.tex (keep a copy of this guide).
\documentclass[11pt]{article}
\usepackage{ailab-preprint}

% ---------- Paper metadata: all front matter is left aligned ----------
\title{A template for AI research\\at Vrije Universiteit Brussel}
\author{First Author\affil{1}, Second Author\affil{1,2}, Third Author\affil{2}}
\affiliations{\affil{1}Artificial Intelligence Lab, Vrije Universiteit Brussel, Belgium\\
  \affil{2}Collaborating institution, City, Country}
\date{Template guide · 17 November 2026}
\preprintabstract{This document is a working template and a short guide for AI research preprints. It demonstrates a single-column layout, collaborative author lists, mathematical statements, pseudocode, numerical results, references, and supplementary material. The opening panel groups the title, authors, and affiliations, followed by a justified abstract; the body uses Libertinus Serif with Libertinus Math. Replace the examples with your research and keep the shared style file unchanged. All numerical results in this guide are illustrative.}
\keywords{research preprints; machine learning; scientific writing}

\begin{document}
\maketitle

\section{Start writing in minutes}
The template keeps formatting in \texttt{ailab-preprint.sty} and research in your
\texttt{.tex} and \texttt{.bib} files. You can write a complete paper without editing
the style. The clean starting point is \texttt{starter.tex}; this guide is
\texttt{main.tex}. Use one as your main document at a time.

\subsection{Open the project in Overleaf}
\begin{enumerate}
  \item Upload the supplied ZIP through \textbf{New project > Existing project (.zip)}.
  \item Open project \textbf{Settings > Compiler}; choose \textbf{LuaLaTeX} and a
        recent TeX Live version (tested locally with TeX Live 2025).
  \item Compile \texttt{main.tex} to read this guide. For your paper, replace its
        contents with \texttt{starter.tex}, then edit the metadata and sections.
  \item Replace \texttt{references.bib} with your bibliography and recompile.
        Overleaf runs BibTeX automatically when needed.
\end{enumerate}
The project needs no shell escape, logo files, or external code highlighting tools.
For the official title font, add the supplied VUB font file as described below. LuaLaTeX loads the requested OpenType fonts directly from TeX Live. Locally, run \texttt{latexmk main.tex}; the included
\texttt{.latexmkrc} selects LuaLaTeX.

\subsection{A restrained research identity}
The palette follows the supplied VUB guide (PDF page 13, printed page 18): blue
\texttt{\#003399}, orange \texttt{\#FF6600}, and gray-blue \texttt{\#E8EBF1}.
The panel uses a lighter tint of gray-blue. Blue carries hierarchy and links;
orange is available as \texttt{VUBOrange} for occasional accents. This is a
research layout informed by the guide. The title uses ITC Avant Garde Gothic Bold
when \texttt{fonts/ITCAvantGardePro-Bold.otf} is supplied. Until then, the preview
uses TeX Gyre Adventor, a substitute, not the official font. Body text and mathematics
retain Libertinus. No full institutional logos are used.

\clearpage % Guide-only page breaks; remove these when writing a paper.
\section{Authors, structure, and everyday writing}
\subsection{Edit the front matter in one place}
Metadata commands go before \verb|\begin{document}|. The superscript helper
\verb|\affil{1,2}| associates authors with any number of institutions. Use explicit
line breaks between affiliations, and let long author lists wrap naturally.
\begin{lstlisting}
\title{A clear and informative research title}
\author{First Author\affil{1}, Second Author\affil{1,2}}
\affiliations{\affil{1}Artificial Intelligence Lab,
  Vrije Universiteit Brussel, Belgium\\
  \affil{2}Partner institution, City, Country}
\preprintabstract{State the problem, approach, key evidence,
  and implications in one or two short paragraphs.}
\keywords{optional; semicolon-separated terms}
\date{17 November 2026} % Optional explicit version date.
\preprintnote{Optional equal-contribution or other author note.}
\end{lstlisting}
Call \verb|\maketitle| once after \verb|\begin{document}|; it prints the justified abstract
below the panel. Do not add a second \texttt{abstract} environment. Omit any
optional command to hide that row. A blank \verb|\date{}| hides the date. For a
large collaboration, the panel can continue onto the next page. Keep the abstract
concise so readers reach the paper quickly. Use \verb|\preprintnote| for
contribution notes in the panel; ordinary body footnotes work normally.\footnote{A body footnote is set in smaller text with a conventional page-local rule.}

\subsection{Organize the argument, then the source}
A typical manuscript introduces the problem and related work, defines the setting,
presents a method and analysis, evaluates it, and discusses limitations.
\citet{saerensPreprint} provides an example combining theory, algorithms,
benchmarks, and a real-world application. Sections can live in separate files:
\begin{lstlisting}
\section{Introduction}\label{sec:intro}
\input{sections/introduction} % File contains prose, not a preamble.
\subsection{Contributions}
\paragraph{Scope.} A short run-in heading for a focused point.
\end{lstlisting}
The title panel takes its cue from the left-aligned opening of
\citet{bessa2026emeans}. The rest of the page stays deliberately simple:
A4, 25\,mm side margins, 11\,pt text, modest paragraph spacing, and centered
page numbers. Section headings use the matching Libertinus Sans family.

\subsection{Optional triangle and official title font}
Use \verb|\usepackage[triangle]{ailab-preprint}| to show the orange VUB triangle
to the left of the title, as in this guide. Omit \texttt{[triangle]} to hide it.
The shape follows the guide's 69.25-degree diagonal (PDF page 32).
To use the official title typeface in Overleaf, upload the licensed bold font as
\texttt{fonts/ITCAvantGardePro-Bold.otf}; it is selected automatically.

\subsection{Cite and cross-reference automatically}
Use \verb|\citet{key}| for a narrative citation and \verb|\citep{key}| for a
parenthetical citation, for example \citep{bessa2026emeans}. Put a
\verb|\label| immediately after a heading or caption and use \verb|\cref| or
\verb|\Cref|: \cref{tab:results,alg:loop} name their object types automatically.
Use \verb|\eqref| for just a parenthesized equation number. Links are blue and
remain legible in monochrome. Long URLs can wrap across lines.

\clearpage
\section{Mathematics, theory, and algorithms}
Write inline mathematics as $x\in\mathbb{R}^{d}$ and use \texttt{equation} or
\texttt{align} for numbered displays. Define paper-specific notation in your
preamble; use \verb|\symbf{x}| for bold mathematical symbols and
\verb|\operatorname| for named operators. Avoid loading legacy font packages
such as \texttt{mathptmx}, \texttt{amssymb}, or \texttt{newtxmath} over
\texttt{unicode-math}.
\begin{align}
  \widehat{L}_n(\theta) &= \frac{1}{n}\sum_{i=1}^{n}
    \ell\bigl(f_\theta(x_i),y_i\bigr), \label{eq:risk}\\
  \theta_{t+1} &= \theta_t-\eta_t\nabla\widehat{L}_n(\theta_t). \label{eq:update}
\end{align}
\begin{lstlisting}
\begin{equation}\label{eq:objective}
  \widehat{L}_n(\theta) = \frac{1}{n}
    \sum_{i=1}^{n}\ell(f_\theta(x_i),y_i)
\end{equation}
Refer to \eqref{eq:objective}, or use \cref{eq:objective}.
\end{lstlisting}

\subsection{Definitions and proofs}
The theorem family includes \texttt{theorem}, \texttt{lemma}, \texttt{proposition},
and \texttt{corollary}, as well as \texttt{conjecture}. For supporting statements,
use \texttt{definition}, \texttt{assumption}, \texttt{remark}, or \texttt{example}.
All share a counter within each section. Use their optional argument for a descriptive name.
\begin{assumption}[Bounded observations]\label{ass:bounded}
The random variables $X_1,\ldots,X_n$ are independent and take values in $[0,1]$.
\end{assumption}
\begin{proposition}[Concentration of the sample mean]\label{prop:bound}
Under \cref{ass:bounded}, let $\bar X_n=n^{-1}\sum_{i=1}^{n}X_i$.
For every $\epsilon>0$,
\[
  \mathbb{P}\!\left(\left|\bar X_n-\mathbb{E}[\bar X_n]\right|\geq\epsilon\right)
  \leq 2\exp(-2n\epsilon^2).
\]
\end{proposition}
\begin{proof}
Apply Hoeffding's inequality to the independent bounded summands. The
upper and lower tails are each at most $\exp(-2n\epsilon^2)$; their union gives
the stated bound.
\end{proof}
Write a statement with \verb|\begin{proposition}[Name]| and end it with
\verb|\end{proposition}|. Follow it with a \texttt{proof} environment; the
end-of-proof symbol is automatic.


\subsection{Pseudocode}
\texttt{algorithm2e} is included. Use \verb|\KwIn|, \verb|\KwOut|,
\verb|\For|, and \verb|\If|; end statements with \verb|\;|.
Do not also load \texttt{algorithm}, \texttt{algorithmic}, or \texttt{algpseudocode}.
\begin{algorithm}[htbp]
\caption{A minimal learning loop (typesetting example)}\label{alg:loop}
\KwIn{Data $D$, initial parameters $\theta_0$, learning rates $(\eta_t)_{t=0}^{T-1}$}
\KwOut{Parameters $\theta_T$}
\For{$t\leftarrow 0$ \KwTo $T-1$}{
  Compute $g_t\leftarrow\nabla\widehat{L}_n(\theta_t)$ on $D$\;
  Update $\theta_{t+1}\leftarrow\theta_t-\eta_t g_t$\;
}
\Return{$\theta_T$}\;
\end{algorithm}

\clearpage
\section{Tables that readers can compare}
Use \texttt{booktabs} rules, descriptive captions above tables, and notes directly
below. Avoid vertical rules, dense grids, and decorative shading. Define the
metric, units, direction of improvement, number of seeds, and meaning of any
uncertainty. Boldface is useful only when its meaning is stated. The values below
are invented to demonstrate formatting; they are not experimental evidence.

\begin{table}[htbp]
\centering
\caption{Illustrative performance over five seeds. Higher accuracy and lower
runtime are better; $\pm$ denotes one sample standard deviation. Bold values
identify the highest mean accuracy, without claiming statistical significance.}
\label{tab:results}
\begin{tabular*}{\linewidth}{@{\extracolsep{\fill}}l
  S[table-format=2.1(2)] S[table-format=3.1(2)] S[table-format=2.1]@{}}
\toprule
Method & {Accuracy (\% $\uparrow$)} & {Runtime (s $\downarrow$)} & {Memory (GB)}\\
\midrule
Baseline & 81.2 +- 0.8 & 125.3 +- 4.2 & 8.4\\
Variant A & 83.7 +- 0.6 & 98.1 +- 3.7 & 9.1\\
Variant B & \multicolumn{1}{c}{\bfseries 85.1 $\pm$ 0.5} & 102.4 +- 2.9 & 9.6\\
\bottomrule
\end{tabular*}
\par\smallskip
{\footnotesize Example setup: identical hardware and evaluation budgets for all methods.
Report confidence intervals instead when appropriate, and describe their construction.\par}
\end{table}

\subsection{Align numbers without manual spacing}
The \texttt{S} column from \texttt{siunitx} aligns decimals and handles uncertainty.
\verb|table-format=3.1(2)| reserves three integer digits, one decimal, and two
uncertainty digits. Write \verb|125.3 +- 4.2|; wrap text headers in braces.
Use \verb|\num{10000}| for \num{10000} and \verb|\qty{12.5}{\second}| for
\qty{12.5}{\second}. Keep units in headers when every row uses the same unit.
\begin{lstlisting}
\begin{table}[htbp]
  \centering
  \caption{Define the metric, uncertainty, and evaluation protocol.}
  \label{tab:comparison}
  \begin{tabular}{@{}l S[table-format=2.1(2)]@{}}
    \toprule
    Method & {Score (\%)}\\
    \midrule
    Baseline & 81.2 +- 0.8\\
    \bottomrule
  \end{tabular}
\end{table}
\end{lstlisting}
For an exceptional bold cell, \verb|\multicolumn{1}{c}{\bfseries ...}|
bypasses numeric parsing; retain the same precision as the other cells.

\subsection{Give descriptive columns room to wrap}
Use \texttt{tabularx} with \verb|\linewidth| and the provided \texttt{Y} column
for prose. Ordinary \texttt{tabular} is better for compact numeric tables.
\begin{table}[htbp]
\caption{Illustrative ablations. Wrapping text preserves readable type at the normal page width.}
\label{tab:ablations}
\begin{tabularx}{\linewidth}{@{}lY S[table-format=2.1]@{}}
\toprule
Configuration & Change from the full method & {Score (\%)}\\
\midrule
Full method & All model components enabled. & 85.1\\
No adaptation & Replace the adaptive schedule with a fixed schedule shared across tasks. & 82.4\\
No regularization & Remove the penalty while keeping the same training budget. & 83.2\\
\bottomrule
\end{tabularx}
\end{table}

\clearpage
\section{Long results and supplementary material}
\subsection{Choose the right table environment}
Use a \texttt{table} float for a table that fits on one page. Use
\texttt{longtable} directly in the document for a multi-page table; it is not
placed inside a \texttt{table} float. Repeat column headings with
\verb|\endfirsthead| and \verb|\endhead|. The appendix below provides an actual
page-spanning example and its source can be copied directly.

For a wide table, first shorten headings, wrap descriptive columns, or split
results into meaningful groups. A local \verb|\small| and slightly reduced
\verb|\tabcolsep| can help. Avoid shrinking an entire table with
\verb|\resizebox|: it can make text and uncertainty unreadable. Put exhaustive
results in an appendix when the main comparison would become hard to scan.

\subsection{Let LaTeX place floats}
Use \verb|[htbp]| and reference every table or algorithm in the text. A float may
move to preserve readable pages; avoid forcing every object to an exact position.
For multi-panel floats, \texttt{subcaption} is already available. Figure design
and figure guidelines are outside this template's current scope.

\subsection{Finish a research manuscript}
After the discussion, include limitations, acknowledgments, and relevant
availability or contribution statements. Acknowledge only funding and resources
that apply to the actual study; do not copy another paper's grants. Keep the
bibliography before the appendix, with one reference list for the whole document.
\begin{lstlisting}
\bibliographystyle{plainnat}
\bibliography{references}
\clearpage
\appendix
\input{appendix} % File begins with \section{...}.
\end{lstlisting}
Call \verb|\appendix| once. Section labels become A, B, and so on; theorem
numbers follow the appendix section. Equations, tables, and algorithms keep
continuous numbering, so every object has a unique number across the PDF.

\subsection{Before sharing the preprint}
Replace sample authors, links, guide prose, and invented numbers. Resolve all
citations and references; inspect equations and tables, and verify author details.
Include the style, bibliography, and every file loaded by \verb|\input| in your
source. The ZIP omits the reference PDF and EPS; neither is needed to compile.

\bibliographystyle{plainnat}
\bibliography{references}

\clearpage
\appendix
\section{A table spanning multiple pages}\label{app:longtable}
This appendix demonstrates repeated headings and continuation notes. The rows
are a synthetic benchmark registry, not a description of a real experiment.
A multi-page table uses the full text width without reducing its font size.
\begingroup
\small
\setlength{\LTleft}{0pt}\setlength{\LTright}{0pt plus 1fill}
\setlength{\LTcapwidth}{\textwidth}
\begin{longtable}{@{}p{.17\textwidth}p{.66\textwidth}>{\raggedleft\arraybackslash}p{\dimexpr.17\textwidth-4\tabcolsep\relax}@{}}
\caption{Illustrative benchmark registry. The rightmost column is an example seed count.}\label{tab:registry}\\
\toprule
Identifier & Evaluation setting & Seeds\\
\midrule
\endfirsthead
\multicolumn{3}{@{}l@{}}{\small\textbf{Table \thetable.} Benchmark registry (continued).}\\[4pt]
\toprule
Identifier & Evaluation setting & Seeds\\
\midrule
\endhead
\midrule
\multicolumn{3}{r@{}}{\footnotesize Continued on the next page.}\\
\endfoot
\bottomrule
\endlastfoot
% Synthetic typesetting examples, not experimental data.
Task 01 & Low-dimensional synthetic task with independent observations. & 5\\[5pt]
Task 02 & High-dimensional synthetic task with correlated features. & 5\\[5pt]
Task 03 & Shifted evaluation distribution with a fixed training budget. & 5\\[5pt]
Task 04 & Sensitivity analysis with a shared hyperparameter configuration. & 5\\[5pt]
Task 05 & Low-dimensional synthetic task with independent observations. & 5\\[5pt]
Task 06 & High-dimensional synthetic task with correlated features. & 5\\[5pt]
Task 07 & Shifted evaluation distribution with a fixed training budget. & 5\\[5pt]
Task 08 & Sensitivity analysis with a shared hyperparameter configuration. & 5\\[5pt]
Task 09 & Low-dimensional synthetic task with independent observations. & 5\\[5pt]
Task 10 & High-dimensional synthetic task with correlated features. & 5\\[5pt]
Task 11 & Shifted evaluation distribution with a fixed training budget. & 5\\[5pt]
Task 12 & Sensitivity analysis with a shared hyperparameter configuration. & 5\\[5pt]
Task 13 & Low-dimensional synthetic task with independent observations. & 5\\[5pt]
Task 14 & High-dimensional synthetic task with correlated features. & 5\\[5pt]
Task 15 & Shifted evaluation distribution with a fixed training budget. & 5\\[5pt]
Task 16 & Sensitivity analysis with a shared hyperparameter configuration. & 5\\[5pt]
Task 17 & Low-dimensional synthetic task with independent observations. & 5\\[5pt]
Task 18 & High-dimensional synthetic task with correlated features. & 5\\[5pt]
Task 19 & Shifted evaluation distribution with a fixed training budget. & 5\\[5pt]
Task 20 & Sensitivity analysis with a shared hyperparameter configuration. & 5\\[5pt]
Task 21 & Low-dimensional synthetic task with independent observations. & 5\\[5pt]
Task 22 & High-dimensional synthetic task with correlated features. & 5\\[5pt]
Task 23 & Shifted evaluation distribution with a fixed training budget. & 5\\[5pt]
Task 24 & Sensitivity analysis with a shared hyperparameter configuration. & 5\\[5pt]
Task 25 & Low-dimensional synthetic task with independent observations. & 5\\[5pt]
Task 26 & High-dimensional synthetic task with correlated features. & 5\\[5pt]
Task 27 & Shifted evaluation distribution with a fixed training budget. & 5\\[5pt]
Task 28 & Sensitivity analysis with a shared hyperparameter configuration. & 5\\[5pt]
Task 29 & Low-dimensional synthetic task with independent observations. & 5\\[5pt]
Task 30 & High-dimensional synthetic task with correlated features. & 5\\[5pt]
Task 31 & Shifted evaluation distribution with a fixed training budget. & 5\\[5pt]
Task 32 & Sensitivity analysis with a shared hyperparameter configuration. & 5\\[5pt]
Task 33 & Low-dimensional synthetic task with independent observations. & 5\\[5pt]
Task 34 & High-dimensional synthetic task with correlated features. & 5\\[5pt]
Task 35 & Shifted evaluation distribution with a fixed training budget. & 5\\[5pt]
Task 36 & Sensitivity analysis with a shared hyperparameter configuration. & 5\\[5pt]
Task 37 & Low-dimensional synthetic task with independent observations. & 5\\[5pt]
Task 38 & High-dimensional synthetic task with correlated features. & 5\\[5pt]
Task 39 & Shifted evaluation distribution with a fixed training budget. & 5\\[5pt]
Task 40 & Sensitivity analysis with a shared hyperparameter configuration. & 5\\[5pt]
Task 41 & Low-dimensional synthetic task with independent observations. & 5\\[5pt]
Task 42 & High-dimensional synthetic task with correlated features. & 5\\[5pt]

\end{longtable}
\endgroup
\section{Extending the paper}\label{app:extensions}
Use additional appendix sections for proofs, hyperparameters, implementation
details, evaluation protocols, and further results. Continue to use descriptive
labels and \verb|\cref|; for example, \cref{tab:registry,prop:bound} remain
unambiguous. For a short appendix, write it directly in your main file. For a
long supplement, use \verb|\input| to keep the source manageable while producing
one PDF with working links and a shared bibliography.
\end{document}
